\documentclass{article}
\usepackage{amsmath}
\usepackage{amssymb}
\usepackage{amsthm}
\usepackage[T1]{fontenc}
\usepackage[normalem]{ulem}

\newtheorem*{conjecture}{Przypuszczenie}
\newtheorem*{theorem}{Twierdzenie}
\newtheorem{lemma}{Lemat}
\renewcommand*{\proofname}{Dowód}


\makeatletter
\newcommand*{\rom}[1]{\expandafter\@slowromancap\romannumeral #1@}
\makeatother


\begin{document}


\section{Przykład 5.}

Pomocniczny dowód do pokazania, że:

\begin{equation*}
    P(X < t) = F_{X}(t^{-})
\end{equation*}

\begin{theorem}
    \begin{equation*}
        (- \infty, t) = \bigcup_{n = 1}^{\infty}(- \infty, t - \frac{1}{n}]        
    \end{equation*}
\end{theorem}

\begin{proof}
    Załóżmy, że \(x \in \bigcup_{n = 1}^{\infty}(- \infty, t - \frac{1}{n}]\). Istnieje zatem indeks \(i \in \mathbb{N}\)
    taki, że \(x \in (- \infty, t - \frac{1}{i}]\). Wiadome jest, że \(t - \frac{1}{n} < t\) dla \(n \in \mathbb{N}\).
    Mamy zatem:
    \begin{equation*}
        - \infty < x < t - \frac{1}{i} < t
    \end{equation*}
    a z tego wynika, że \(x \in (- \infty, t)\).

    W drugą stronę, niech \(x \in (- \infty, t)\). Załóżmy dla uzyskania sprzeczności, że \(x \not\in \bigcup_{n = 1}^{\infty}(- \infty, t - \frac{1}{n}] \).
    Wtedy:
    \begin{equation*}
        x \in \bigcap_{n = 1}^{\infty}(t - \frac{1}{n}, + \infty)
    \end{equation*}
    Połóżmy \(x_n = x + \frac{1}{n} \). Taki ciąg ma granicę \(\lim_{n \to \infty} x_n = x\). Zauważmy, że \(x_n \in (t - \frac{1}{n}, + \infty)\) dla dowolnego \(n \in \mathbb{N}\).
    Mamy wobec tego:
    \begin{equation*}
        t - \frac{1}{n} < x_n < \infty
    \end{equation*}
    Mamy również:
    \begin{equation*}
        t - \frac{1}{n} < x_n <  t + \frac{1}{n}
    \end{equation*}
    \sout{I przechodząc do granicy otrzymujemy} Źle, nierówności ostre się nie przenoszą na granice!:
    \begin{equation*}
        t < x < t
    \end{equation*}
    czyli \(x \in \emptyset\), ale z założenia \(x \in (- \infty, t)\). Sprzeczność.
\end{proof}

\subsection*{Granica jednostronne, suprema oraz końcówka wzoru}

\emph{Spróbować pokazac, ze jak F jest monotoniczna na jakims przedziale,
to w kazdym punkcie tego przedzialu taka granica istnieje}.

To znaczy, udowodnimy, że jak \(F\) rosnie (w ogólności jest po prostu monotoniczna), to:

\begin{equation*}
    \lim_{s\to t^{-}} F(s) = \sup_{s \in (-\infty, t)} F(s)
\end{equation*}

Ustalmy \(\epsilon > 0\). Pokażemy narazie szkic na objawienie. Musimy znaleźć \(\delta > 0\) zachowującą się w taki sposób:

\begin{equation*}
    \sup F(s) - \epsilon < F(t - \delta)
\end{equation*}

Pomocnicze będzie to (operacje algebraiczne i fakt, że \(s < t\)):

\begin{equation*}
    0 < t - s < \delta \implies - \delta < t - s - \delta < 0 \implies s - \delta < t - \delta < s < t
\end{equation*}

Zaczynamy. Dobieramy sobie \(s'\) większe od \(t - \delta\) i mniejsze niż \(t\) (podkładka pod monotoniczność) i zajdzie:

\begin{equation*}
    \sup F(s) - \epsilon < F(t - \delta) \leq F(s') \leq \sup F(s)
\end{equation*}

Z tego mamy (dostrzeż ogon/koniec z prawej strony)

\begin{equation*}
    \sup F(s) - \epsilon < F(s') \leq \sup F(s) < \sup F(s) + \epsilon
\end{equation*}
Własność wartości bezwzględnej:
\begin{equation*}
    |F(s') - \sup F(s)| < \epsilon
\end{equation*}
Niech \(M = \sup \{ F(s) : s \in (- \infty, t) \} = \sup_{s \in (-\infty, t)} F(s)\). Skoro \(M\) jest supremum to
korzystamy z tego, że dla każdego \(\epsilon > 0\) istnieje \(F(z) \in \{ F(s) : s \in (- \infty, t) \}\) taki,
że \(M - \epsilon < F(z)\). Podstawiamy zatem za \(z = t - \delta\). Szukana \(\delta\) to \(t - z\).


\end{document}
